% Einheit 3 von 3 – Einheit 3 (bank/normalverteilung-und-sigma-regeln/e3.jsonl, zusammenbau v0.1)
%% TODO Prüfungswort Einheit 3: nicht in der Marken-Zeile
%% TODO Zeitmarke Einheit 3: Marken-Zeile nicht lesbar
%% TODO Zweigzeile Teil 1: was hier gelernt wird (Satz in Schülersprache aus dem Einheitstitel) – Einheit 3
%% TODO Einheitentitel 3 nicht in der Mappe lesbar
\einheitenkopf[e3]{Einheit 3 von 3 · Einheit 3}
\setcounter{aufgabe}{11}

%% TODO Titel: Ich-kann-Satz fehlt in der Bank (Platzhalter: Kettenname) – Nr. 12
%% TODO Anweisung: ein Satz über der ersten Teilaufgabe fehlt in der Bank – Nr. 12
\begin{aufgabe}{Umkehraufgaben und Argumente}
%% TODO \rechenplatz{n}: Schrittzahl des Grundfalls fehlt in der Bank (nur bei fester Schreibform, 2.3 b)
\begin{teile}
\teil $\sigma = 2$ ist bekannt; $P(X > 10)$ soll $5\,\%$ betragen. Wie groß muss $\mu$ sein? \\ \kreuz{vorwärts: eine Wahrscheinlichkeit ist gesucht} \\ \kreuz{Umkehraufgabe: ein Parameter oder eine Grenze ist gesucht}
\teil Die Dichte von $X$ ist $f(x) = \frac{1}{3\sqrt{2\pi}} \, \mathrm{e}^{-\frac{1}{2}\left(\frac{x - 40}{3}\right)^2}$. $\mu$ und $\sigma$? \hfill (Abitur 2025 LK) \\ μ = \leerfeld , σ = \leerfeld
\teil Die Wartezeit $W$ an einem Schalter in Minuten ist normalverteilt mit $\sigma = 2$. Nur $5\,\%$ der Kunden warten länger als $10$ Minuten. Erwartungswert? Dann $P(W < 5)$? \hfill (Abitur 2022 LK)
\teil Reifen A: Laufleistung normalverteilt mit $\mu = 45\,000$ km, $\sigma = 5\,000$ km. Bestimme die Laufleistung $c$, die $90\,\%$ der Reifen A übertreffen. Reifen B: $\mu = 42\,000$ km, $\sigma = 3\,000$ km. Übertreffen auch $90\,\%$ der Reifen B diese Laufleistung? \hfill (Abitur 2024 LK)
\teil Eine Anlage füllt Dosen; die Füllmenge $X$ in g ist normalverteilt mit $\mu = 500$ und unbekanntem $\sigma$. Bedingung I: $P(X \le 494) \le 5\,\%$. Bedingung II: $P(X \le 488) \le 0{,}1\,\%$. Begründe: Erfüllt die Anlage Bedingung I, so auch Bedingung II. \hfill (Abitur 2025 LK)
\end{teile}
\end{aufgabe}

%% TODO Titel: Ich-kann-Satz fehlt in der Bank (Platzhalter: Kettenname) – Nr. 13
%% TODO Anweisung: ein Satz über der ersten Teilaufgabe fehlt in der Bank – Nr. 13
\begin{aufgabe}{Standardabweichung einer Normalverteilung aus einer vorgegebenen Wahrscheinlichkeit bei festem Erwartungswert ermitteln}
\begin{teile}
\teil Die Länge $X$ einer Schraube in mm ist normalverteilt mit $\mu = 40$ und $\sigma = 0{,}5$; kürzer als $39$ mm ist Ausschuss. Nach einer Verbesserung ist die Wahrscheinlichkeit dafür halbiert, $\mu$ bleibt gleich. Neue Standardabweichung? \hfill (Abitur 2017 LK) \\ σ = \leerfeld
\end{teile}
\end{aufgabe}

%% TODO Titel: Ich-kann-Satz fehlt in der Bank (Platzhalter: Kettenname) – Nr. 14
%% TODO Anweisung: ein Satz über der ersten Teilaufgabe fehlt in der Bank – Nr. 14
\begin{aufgabe}{Umkehraufgaben und Argumente – Fehler finden, Begründen}
\begin{teile}
\teil Ich finde den Fehler: Die Dichte von $X$ ist $f(x) = \frac{1}{\sqrt{18\pi}} \, \mathrm{e}^{-\frac{(x - 30)^2}{18}}$. Jana: „$\mu = 30$, $\sigma = 18$.“
\teil Begründe: Am Graphen der Verteilungsfunktion liegt $\mu$ an der Stelle mit dem Funktionswert $\frac{1}{2}$.
\end{teile}
\end{aufgabe}
